\chapter{Desarrollo de la investigación}



\section{Concepto A en el Desarrollo de la Investigación con el input}

\lipsum[1-1]

\dots se detallan en el Apéndice~\ref{AppendA}.

\section{Teorías C  en el Desarrollo de la Investigación con el input}

\lipsum[1-1]

\dots en la Tabla \ref{table:movimientos_basicos}.

\begin{table}[H]
\centering
\begin{tabular}{| c | c | c|} 
 \hline
 \textbf{Modo} & \textbf{Articulaciones} & \textbf{Ruedas}\\ [0.3ex] 
 \hline \hline
 Avanza & $J_1=J_4=512$, $J_2=J_3=-512$ & $W_1=W_3=-v$, $W_2=W_4=v$  \\ 
 \hline
 Retrocede & $J_1=J_4=512$, $J_2=J_3=-512$ & $W_1=W_3=v$, $W_2=W_4=-v$  \\
 \hline \hline 
 Giro CCW & $J_1=J_2=J_3=J_4=0$ & $W_1=W_2=W_3=W_4=v$ \\
 \hline 
 Giro CW & $J_1=J_2=J_3=J_4=0$ & $W_1=W_2=W_3=W_4=-v$  \\ %[1ex] 
 \hline \hline 
 Derecha & $J_1=J_4=-512$, $J_2=J_3=512$ & $W_1=W_2=-v$, $W_3=W_4=v$  \\ %[1ex] 
 \hline
 Izquierda & $J_1=J_4=-512$, $J_2=J_3=512$ & $W_1=W_2=v$, $W_3=W_4=-v$  \\
 \hline
\end{tabular}
\caption[Movimientos Básicos Programados]
	{Movimientos Básicos Programados. $J_i$ es el valor de la posición para el servomotor que controla la articulación de la i-ésima extremidad, y donde $W_i$ es el valor de la velocidad de la rueda de la i-ésima extremidad.  Fuente: Elaboración propia.}
\label{table:movimientos_basicos}
\end{table}

\section{Metodología E  en el Desarrollo de la Investigación con el input}
\lipsum[1-1]

\begin{lstlisting}[language=Python, caption=Ejemplo en Python, label={lst:alg1}, numbers=left]
import numpy as np
    
def sum(x,y):
    return x + y
\end{lstlisting}

\lipsum[1-1]

\dots el Algoritmo \ref{lst:alg1}.

%\afterpage{\blankpage}